\input{../tex-inputs/latex-leseansicht-vorspann}

\section[Arthur Schnitzler an Gustav Schwarzkopf, 27. 7. 1899]{L04043 Arthur Schnitzler an Gustav Schwarzkopf, 27. 7. 1899}
\nopagebreak\mylabel{L04043v}
\rehead{ }\normalsize\beginnumbering\briefempfaengerindex{Schwarzkopf, Gustav@\textsc{Schwarzkopf, Gustav}!zzzSchnitzler, Arthur@\emph{von Arthur Schnitzler}!1899-07-272@{27. 7. 1899}|(be}
\toendnotes[C]{\smallbreak\pagebreak[2]}
\correspDesc{Versand  durch Arthur Schnitzler am 27. 7. 1899 in Velden am Wörthersee
\newline{}Erhalt  durch Gustav Schwarzkopf im Zeitraum [28. 7. 1899 – 29. 7. 1899] in Wien}\toendnotes[C]{\smallbreak}
\Standort{CUL, Schnitzler, B 96.}
\physDesc{Brief, 1 Blatt, 3 Seiten, 1467 Zeichen
\newline{}Handschrift: Bleistift, deutsche Kurrent}\toendnotes[C]{\smallbreak}
\pstart
           \raggedleft{}{\pb}\textsc{Velden, Pens. Pundschu\oindex{Pension Pundschu@\textbf{Pension Pundschu}, \emph{Hotel}|pw}}\pend
           
\pstart
           \raggedleft{}Donnerſtg\pend
           \vspace{0.5em}
\pstart
           lieber Guſtav, noch immer bin ich hier, wie Sie ſehen; heute reiſen
                  Mama\pwindex{Schnitzler, Louise 8.\,7.\,1840 Kőszeg – 9.\,9.\,1911 Wien@\textsc{Schnitzler, Louise} (8.\,7.\,1840 Kőszeg – 9.\,9.\,1911 Wien)|pwv} u Giſela\pwindex{Hajek, Gisela 20.\,12.\,1867 Wien – 3.\,2.\,1953 Cambridge@\textsc{Hajek, Gisela} (20.\,12.\,1867 Wien – 3.\,2.\,1953 Cambridge)|pw} nach Toblach\oindex{Toblach@\textbf{Toblach}, \emph{Verwaltungsgebiet}|pw}; ich
               und \textsc{Wasserm.\pwindex{Wassermann, Jakob 10.\,3.\,1873 Fürth – 1.\,1.\,1934 Altaussee@\textsc{Wassermann, Jakob} (10.\,3.\,1873 Fürth – 1.\,1.\,1934 Altaussee), \emph{Schriftsteller}|pw}}{ }\label{K_L04043-1v}\edtext{morgen nach Villach\oindex{Villach@\textbf{Villach}, \emph{Verwaltungsgebiet}|pw}}{\lemma{\textnormal{\emph{morgen nach Villach}}}\Cendnote{\textnormal{Vgl. A. S.: \emph{Tagebuch}, 28. 7. 1899.}}}\label{K_L04043-1}, ſind So{\geminationn}tag bei Richard\pwindex{Beer-Hofmann, Richard 11.\,7.\,1866 Wien – 26.\,9.\,1945 New York City@\textsc{Beer-Hofmann, Richard} (11.\,7.\,1866 Wien – 26.\,9.\,1945 New York City), \emph{Schriftsteller}|pw}; der \label{K_L04043-2v}\edtext{neulich bei uns}{\lemma{\textnormal{\emph{neulich bei uns}}}\Cendnote{\textnormal{A. S.: \emph{Tagebuch}, 24. 7. 1899.}}}\label{K_L04043-2} war und in viel
               beſſrer Sti{\geminationm}ung als ich vermuthet. Wir wollen alſo
               wirklich die berühmte Fußpartie machen, Richard\pwindex{Beer-Hofmann, Richard 11.\,7.\,1866 Wien – 26.\,9.\,1945 New York City@\textsc{Beer-Hofmann, Richard} (11.\,7.\,1866 Wien – 26.\,9.\,1945 New York City), \emph{Schriftsteller}|pw}, Waſſerm\pwindex{Wassermann, Jakob 10.\,3.\,1873 Fürth – 1.\,1.\,1934 Altaussee@\textsc{Wassermann, Jakob} (10.\,3.\,1873 Fürth – 1.\,1.\,1934 Altaussee), \emph{Schriftsteller}|pw}, ich – am
               Ende ſchließt{ }ſich auch Robert Hirschf.\pwindex{Hirschfeld, Robert 17.\,9.\,1857 Žďár nad Sázavou – 2.\,4.\,1914 Salzburg@\textsc{Hirschfeld, Robert} (17.\,9.\,1857 Žďár nad Sázavou – 2.\,4.\,1914 Salzburg), \emph{Journalist, Musikkritiker}|pw} an –
               und we{\geminationn} Sie mitgehen wollten, wär es unter Ihren vielen
               guten Ideen nicht die ſchlechteſte. Überlegen Sie ſichs doch. Wir brechen etwa
                  5., od. 6. Auguſt von Toblach\oindex{Toblach@\textbf{Toblach}, \emph{Verwaltungsgebiet}|pw} auf, das ganze ni{\geminationm}t etwa 10 Tage {\pb}in Anſpruch, endet in Bozen\oindex{Bozen@\textbf{Bozen}, \emph{Hauptstadt}|pw}, von wo \introOben{}(oder Trient\oindex{Trient@\textbf{Trient}|pw})\introOben{} ich etwa Mitte Auguſt nach Iſchl\oindex{Bad Ischl@\textbf{Bad Ischl}|pw} fahre. Schreiben Sie mir nach \textsc{Toblach\oindex{Toblach@\textbf{Toblach}, \emph{Verwaltungsgebiet}|pw}}, Südbahnhotel\oindex{Südbahnhotel [Toblach]@\textbf{Südbahnhotel [Toblach]}, \emph{Hotel}|pw}. \label{K_L04043-3v}\edtext{Montag bin ich ſchon dort}{\lemma{\textnormal{\emph{Montag … dort}}}\Cendnote{\textnormal{A. S.: \emph{Wiener Schnitzler}, 31. 7. 1899.}}}\label{K_L04043-3}.
               Ich{ }ſetze nichts weiter hinzu – hoffentlich iſt es wirklich unnöthig. Sind Sie
               abſolut gegen die Fußpartie –{ }ſo kommen Sie Mitte Auguſt nach Iſchl\oindex{Bad Ischl@\textbf{Bad Ischl}|pw}, wo nicht nur ich, – ſondern auch meine
               kleinen Neffen\pwindex{Schnitzler, Hans 11.\,7.\,1895 Wien – 26.\,3.\,1967 Chicago@\textsc{Schnitzler, Hans} (11.\,7.\,1895 Wien – 26.\,3.\,1967 Chicago), \emph{Chirurg}|pwv}\pwindex{Schnitzler, Karl 11.\,10.\,1896 Wien – 12.\,9.\,1981 Sydney@\textsc{Schnitzler, Karl} (11.\,10.\,1896 Wien – 12.\,9.\,1981 Sydney), \emph{Geschäftsführer}|pwv} ſind. Aber das beſte wäre, Sie nähmen Südtirol\oindex{Südtirol@\textbf{Südtirol}, \emph{Verwaltungsgebiet}|pw} und Iſchl\oindex{Bad Ischl@\textbf{Bad Ischl}|pw} zusa{\geminationm}en. – Ich habe hier wieder \label{K_L04043-4v}\edtext{zu arbeiten\pwindex{Schnitzler, Arthur 15.\,5.\,1862 Wien – 21.\,10.\,1931 ebd.@\textsc{Schnitzler, Arthur} (15.\,5.\,1862 Wien – 21.\,10.\,1931 ebd.), \emph{Schriftsteller, Mediziner}!Schleier der Beatrice. Schauspiel in fünf Akten@\strich\emph{Der Schleier der Beatrice. Schauspiel in fünf Akten}|pwv} angefangen}{\lemma{\textnormal{\emph{zu arbeiten angefangen}}}\Cendnote{\textnormal{Vgl. A. S.: \emph{Tagebuch}, 23. 7. 1899.}}}\label{K_L04043-4}; mit
               mäßigem Talent. {\pb}Auch die Partie ſoll
               theilweise zum Arbeiten verwendet werden;– die Tage ſind ſo lang. Das Wetter hier war
               faſt durchwegs ſchön; das tägliche Bad ſehr angenehm; einſame Fahrten um den See\oindex{Wörthersee@\textbf{Wörthersee}, \emph{See}|pw} von großer Traurigkeit; die Abende manchmal
               kaum erträglich. Nachmittag hab ich meiſtens geſchrieben. – Die zweite Hälfte Auguſt
               denk ich in Iſchl\oindex{Bad Ischl@\textbf{Bad Ischl}|pw} zu verbringen; im
                  September erſt durch Thüringen\oindex{Thüringen@\textbf{Thüringen}, \emph{Land}|pw} zu
               radeln. – Bitte kommen Sie!\pend
           \pstart Von Herzen Ihr \spacefill\mbox{ArthSch.}\pend{}\selectlanguage{ngerman}\endnumbering\briefempfaengerindex{Schwarzkopf, Gustav@\textsc{Schwarzkopf, Gustav}!zzzSchnitzler, Arthur@\emph{von Arthur Schnitzler}!1899-07-272@{27. 7. 1899}|)be}\mylabel{L04043h}
\begin{anhang}
\end{anhang}\newcommand{\dateiname}{L04043}\newcommand{\titel}{Arthur Schnitzler an Gustav Schwarzkopf, 27. 7. 1899}\newcommand{\editorInnen}{Herausgegeben von Jahnke, SelmaMüller, Martin Anton}\input{../tex-inputs/latex-leseansicht-abspann}
